\chapter{Anatomía del \emph{ingest}}
\label{cap:ingest-anatomy}

\begin{intuicion}
``Ingesta'' es el camino \emph{log $\to$ grafo}. Es el componente
con menos magia algorítmica del libro pero el de más impacto en
escalabilidad: si el ingestor lee el archivo entero a RAM, una
sesión de 28~GB es inviable; si lo procesa en \emph{batches} de 50K
filas el peak RAM queda en $\sim$10~MB sin importar el tamaño del
archivo. Tres componentes orquestan el flujo: \texttt{LogSource}
(``decime el próximo chunk de bytes''), \texttt{LogParser}
(``convertí estos bytes en triples''), y \texttt{GraphWriter}
(``empujá estos triples al grafo''). Conectados por canales async,
funcionan como un pipeline donde cada etapa procesa lo que la
anterior produjo, sin acumular.
\end{intuicion}

\section{Diagrama del pipeline}

\begin{figure}[H]
\centering
\begin{tikzpicture}[
    node distance=0.5cm,
    every node/.style={font=\footnotesize, align=center},
    box/.style={draw=accent, thick, rounded corners=1.5mm,
                minimum height=1.2cm, minimum width=2.6cm,
                fill=accentlight}
]
\node[box] (src) {\textbf{LogSource}\\\scriptsize Sentinel, Defender,\\\scriptsize Data Lake, EVTX};
\node[box, right=of src] (poll) {\textbf{SourcePoller}\\\scriptsize \emph{async pump}};
\node[box, right=of poll] (parse) {\textbf{LogParser}\\\scriptsize Sysmon, Sentinel,\\\scriptsize Generic, Forti};
\node[box, right=of parse] (write) {\textbf{GraphWriter}\\\scriptsize add\_entity\\\scriptsize add\_relation};
\node[box, right=of write] (store) {\textbf{streaming}\\\scriptsize auto-spill mmap\\\scriptsize canonical store};

\foreach \i/\j in {src/poll, poll/parse, parse/write, write/store} {
    \draw[->, >=Latex, accent, thick] (\i) -- (\j);
}

\node[below=0.1cm of poll, font=\scriptsize\itshape, accent] {watermarks};
\node[below=0.1cm of parse, font=\scriptsize\itshape, accent] {triples};
\node[below=0.1cm of write, font=\scriptsize\itshape, accent] {single-write};
\end{tikzpicture}
\caption{Pipeline de ingesta. Cada caja es un componente con su
propio módulo en \texttt{platform/sources}, \texttt{platform/parsers}
o \texttt{core/graph-engine}. Las flechas son canales async
(\texttt{tokio::sync::mpsc}).}
\label{fig:ingest-pipeline}
\end{figure}

\section{Etapa 1: \texttt{LogSource}}

\begin{lstlisting}[language=Rust]
pub trait LogSource: Send + Sync {
    type Chunk;
    async fn next_chunk(&mut self) -> Option<Self::Chunk>;
    fn watermark(&self) -> Option<i64>;  // último ts visto
}
\end{lstlisting}

\noindent
Una \texttt{LogSource} es un cursor sobre una fuente concreta. El
\texttt{SourcePoller} ---el componente que orquesta el pipeline---
le pide chunks una y otra vez hasta que devuelve \texttt{None}.

\subsection{Implementaciones}

\begin{itemize}
    \item \texttt{SentinelSource}: query KQL contra Microsoft
    Sentinel. Pagina hasta que la API devuelve un cursor vacío.
    \item \texttt{DefenderXdrSource}: similar, para Microsoft
    Defender XDR.
    \item \texttt{DataLakeSource}: lee archivos Parquet desde un
    bucket S3-compat. Itera por partición temporal.
    \item \texttt{EvtxSource}: lee archivos \texttt{.evtx} binarios
    de Windows. Usa el crate \texttt{evtx} para parsear el formato
    binario.
    \item \texttt{NdjsonFileSource}: archivos NDJSON locales (uso
    típico para datasets de prueba).
\end{itemize}

\subsection{Watermarks}

Algunas fuentes tienen ordering temporal natural (Sentinel pagina
por timestamp); otras no (Data Lake puede entregar particiones en
cualquier orden). \texttt{watermark()} devuelve el menor timestamp
de eventos \emph{garantizado-ya-vistos}: el ingestor lo usa para
saber hasta qué momento puede hacer \emph{commit} de aristas en
modo \emph{streaming}.

\section{Etapa 2: \texttt{LogParser}}

\begin{lstlisting}[language=Rust]
pub trait LogParser: Send + Sync {
    type Chunk;
    fn parse(&self, chunk: Self::Chunk)
        -> Result<Vec<ParsedTriple>, ParseError>;
}
\end{lstlisting}

\noindent
El \texttt{LogParser} convierte un chunk crudo (bytes, NDJSON, EVTX
binario) en una secuencia de \texttt{ParsedTriple} (capítulo
\ref{cap:entity-relation}). Cada parser concreto vive en
\filepath{platform/parsers/}:

\begin{itemize}
    \item \texttt{SysmonNdjsonParser}: NDJSON de Sysmon (Event ID
    1, 3, 7, 8, 11, etc.). Cada Event ID se mapea a uno o varios
    triples.
    \item \texttt{SentinelKqlParser}: rows JSON del schema canonical
    de Sentinel.
    \item \texttt{EvtxBinaryParser}: el XML EVTX binario crudo.
    \item \texttt{GenericCsvParser}: CSV con \emph{field config}
    explícito (qué columna es source, qué columna es timestamp,
    etc.). Usado para datasets ad-hoc.
    \item \texttt{CognitoParser}, \texttt{FortianalyzerParser},
    \texttt{IisParser}: formatos verticales específicos.
\end{itemize}

\subsection{El parser es \emph{event-to-N-triples}}

Un evento de log puede generar 0, 1, o varios triples. El parser de
Sysmon Event ID 1 (Process Create) genera 2 triples (host-hosts-process,
parent-spawn-process). El parser de un Event ID irrelevante (un
debug log) genera 0. El parser de un evento ``compuesto''
(autenticación con MFA) puede generar 3 (user-auth-host,
host-mfa-device, user-owns-device).

\section{Etapa 3: \texttt{GraphWriter}}

\begin{lstlisting}[language=Rust]
pub struct GraphWriter {
    graph: Arc<RwLock<GraphHunter>>,
}

impl GraphWriter {
    pub fn apply(&self, triple: ParsedTriple) {
        let mut g = self.graph.write();
        g.add_entity(Entity {
            id: triple.source_id.clone(),
            entity_type: triple.source_type,
            metadata: HashMap::new(),
        });
        g.add_entity(Entity {
            id: triple.dest_id.clone(),
            entity_type: triple.dest_type,
            metadata: HashMap::new(),
        });
        g.add_relation(Relation {
            source_id: triple.source_id,
            dest_id: triple.dest_id,
            rel_type: triple.rel_type,
            timestamp: triple.timestamp,
            metadata: triple.metadata,
        });
    }
}
\end{lstlisting}

\noindent
\texttt{add\_entity} es idempotente (deduplicado por el interner,
capítulo~\ref{cap:interner}); \texttt{add\_relation} es append +
\emph{streaming mirror} (capítulo~\ref{cap:streaming-mirror}).

\section{El \texttt{SourcePoller}}

\begin{lstlisting}[language=Rust]
pub async fn run<S, P>(
    mut source: S,
    parser: Arc<P>,
    writer: Arc<GraphWriter>,
) -> Result<IngestStats, Error>
where
    S: LogSource,
    P: LogParser<Chunk = S::Chunk>,
{
    let mut stats = IngestStats::default();
    while let Some(chunk) = source.next_chunk().await {
        let triples = parser.parse(chunk)?;
        stats.events += 1;
        stats.triples += triples.len();
        for triple in triples {
            writer.apply(triple);
        }
    }
    Ok(stats)
}
\end{lstlisting}

\noindent
El loop es deliberadamente simple: una fila a la vez, sin batching
explícito. La razón: con \texttt{add\_relation} pesando $\sim$1
$\mu$s (interner lookup + push al store + push al streaming), el
overhead de batching no compensa la complejidad.

\subsection{Concurrencia: ¿múltiples writers?}

GraphHunter mantiene el ingestor \emph{single-threaded}. Razón: el
\texttt{GraphHunter} interno tiene varias estructuras (interner,
edge store, streaming, índices invalidations) que serían difíciles
de coordinar bajo varios writers. El bottleneck efectivo no es CPU
del ingestor sino I/O de la fuente; un solo thread alcanza para
saturar Sentinel KQL o un disco SSD.

Cuando hace falta acelerar (datasets multi-fuente), se corren
\emph{varios SourcePollers en paralelo} ---uno por fuente--- pero
todos llaman al mismo \texttt{GraphWriter} bajo lock. El lock en
\texttt{add\_relation} es un \texttt{RwLock} en write mode; los
matchers concurrentes usan read mode y no se bloquean entre sí.

\begin{nota}
El streaming backend (capítulo~\ref{cap:streaming-mirror}) está
diseñado para soportar mejor la concurrencia: cada vértice tiene su
propia \texttt{NeighborList} con sincronización fina. La migración
del matcher para usar streaming como source of truth (stage 2c) abre
la puerta a varios writers concurrentes; ese cambio está en
diseño futuro (capítulo~\ref{cap:streaming-canonical}).
\end{nota}

\section{Memoria del ingestor: peak vs total}

Una propiedad importante: el \emph{peak} RAM del ingestor es del
orden del \emph{batch size}, no del tamaño del archivo.

\begin{itemize}
    \item Chunk típico: 50.000 NDJSON rows, $\sim$10~MB de bytes
    crudos.
    \item Triples generados: $\sim$25.000 (algunos eventos no
    generan triples, otros generan varios).
    \item Memoria viva por chunk: $\sim$10~MB de bytes + $\sim$5~MB
    de \texttt{ParsedTriple}s. Se libera cuando el chunk se
    consumió.
\end{itemize}

\noindent
La memoria \emph{total} consumida (lo que se acumula en el grafo)
sí crece con el tamaño del log: 28~GB de EVTX produce $\sim$200~M
aristas $\Rightarrow$ $\sim$6.4~GB de \texttt{CompactRelation}.
Pero \emph{ese} crecimiento es lineal y predecible; el auto-spill
del \emph{streaming backend} (capítulo~\ref{cap:streaming-mirror})
lo absorbe cuando excede el budget.

\section{Provenance}

Cada triple lleva implícitamente la información de \emph{de qué
fuente vino}: el \texttt{SourcePoller} setea
\texttt{current\_dataset\_tag} en el grafo antes de empezar, y
\texttt{add\_relation} lo copia a \texttt{CompactRelation::dataset\_tag}
(capítulo~\ref{cap:compact-relation}). Para sesiones multi-fuente,
los matchers pueden filtrar por \texttt{dataset\_tag} para
restringir hunts a un subset.

\section{Resumen}

\begin{itemize}
    \item Pipeline tres-etapas: \texttt{LogSource} (chunks),
    \texttt{LogParser} (triples), \texttt{GraphWriter}
    (entities/relations).
    \item Async via \texttt{tokio::sync::mpsc}; canales conectan las
    etapas.
    \item Peak RAM = batch size, no archivo. Total RAM = aristas
    acumuladas (con spill cuando excede budget).
    \item Single-threaded por simplicidad; bottleneck es I/O de la
    fuente, no CPU.
    \item Provenance vía \texttt{dataset\_tag} en cada arista.
\end{itemize}

\begin{ejercicios}
\begin{enumerate}
    \item \starred{} Si una fuente Sentinel devuelve 1.000 eventos
    por segundo y cada uno genera 1.5 triples promedio, ¿cuántas
    aristas/segundo agregás al grafo? Si la sesión dura 8 horas,
    ¿cuántas aristas en total?

    \item Dibujá el peak RAM del ingestor a lo largo del tiempo
    para una sesión de 28~GB. ¿Tiene picos? ¿Por qué se estabiliza
    en $\sim$10~MB del chunk?

    \item \starred{} El loop del \texttt{SourcePoller} es
    \texttt{while let Some(chunk) = source.next\_chunk().await}.
    ¿Qué pasa si la fuente nunca devuelve \texttt{None} (un
    streaming infinito de Sentinel pulling)? Diseñá una
    \emph{stop condition} adecuada.

    \item Imaginá que querés correr 10 \texttt{SourcePoller}s en
    paralelo (10 fuentes distintas). ¿Cómo afecta al
    \texttt{RwLock} del \texttt{GraphHunter}? ¿Qué cuello de
    botella se vuelve dominante?

    \item ¿Por qué \texttt{add\_entity} es idempotente y
    \texttt{add\_relation} no? Discutí qué cambia para el ingestor
    si lo corriera con \texttt{add\_relation} idempotente
    (deduplicando por terna).
\end{enumerate}
\end{ejercicios}
